\chapter{Conclusions and Future Directions}
\label{chap:conclusion}

Chapter~\ref{chap:learned_control} closed with the observation that the two studies of that chapter
complete the progression announced in Chapter~1: perception locates a component before either arm
can act on it, a second arm cannot be added safely before it is decided which arm acts, concurrent
motion cannot be trusted without a continuous safety signal, that signal is only useful once a
planner consumes it, and only with a working classical planner in place does it become possible to
identify precisely where a hand-designed control law runs out of headroom and a learned one is
warranted. This chapter closes the thesis itself. Section~\ref{sec:concl_bygap} revisits each of
the seven research gaps identified in Chapter~\ref{chap:lit_review}, restating the contribution the
corresponding study makes against it, and then, for each study, stating the limitations and the
future directions the study's own authors identified, in the terms the original work itself used.
Section~\ref{sec:concl_remarks} closes with brief concluding remarks on the thesis as a whole.

%% ============================================================
\section{Contributions, Limitations, and Future Directions by Research Gap}\label{sec:concl_bygap}
%% ============================================================

The eight studies of this thesis were organized in Chapter~1 around seven research gaps, G1 through
G7, each closed by one or more of the studies. What follows revisits each gap in turn: the
contribution restates, briefly, what Chapters~\ref{chap:vision_pipelines}
through~\ref{chap:learned_control} already established in full; the limitations and future
directions are drawn directly from the discussion each constituent study gives of its own scope,
rather than from any assessment introduced in this chapter.

%% ------------------------------------------------------------
\subsection{G1: Perception for Small, Low-Feature Components}\label{sec:concl_g1}
%% ------------------------------------------------------------

\begin{description}
    \item[Contribution.] The first study localizes overlapping bolts from RGB-D point clouds using
    $k$-means segmentation to restrict the search volume and Oriented Chamfer matching for pose
    estimation, with colour-based segmentation and contour analysis for size estimation, achieving
    bolt sorting from a cluttered bin without misclassification in Gazebo
    simulation~\citep{sk2024bolthandling}. The second study replaces classical segmentation with a
    custom-trained DeepLabV3+ model, reaching a mean Intersection over Union of 0.95 against
    PPLiteSeg and OCRNet baselines on a custom dataset of new and aged bolts, and fuses the resulting
    mask with RealSense D415 depth to localize and sort bolts by size on hardware, again without
    misclassification~\citep{sk2024boltsorting}.

    \item[Limitations.] For the first study, position and size accuracy are bounded by the
    point-cloud step size used during Chamfer matching, by the decision to ignore template scale
    during matching, and by camera resolution, higher resolution reduced size-estimation error but
    increased computation time; validation was limited to Gazebo simulation with bolts assumed
    static and handled with a stock gripper~\citep{sk2024bolthandling}. For the second study, bolts
    were likewise assumed static and handled with the in-built gripper, the pipeline was validated
    on bolts presented singly rather than in overlapping or occluded configurations or in motion, and
    only bolts, not nuts, were sorted~\citep{sk2024boltsorting}.

    \item[Future directions.] The first study's stated future work is to optimize computational
    efficiency through GPU acceleration and geometric descriptors, to expand the system to a more
    diverse set of components, and to integrate the perception pipeline with a dual-arm robotic
    system~\citep{sk2024bolthandling}, a direction the third study of this thesis subsequently takes
    up directly. The second study's stated future work is to handle bolts in overlapping or occluded
    configurations and in motion, to reduce sorting time using a customized gripper, and to extend
    the pipeline to include nuts for nut-bolt assembly using a dual-arm
    manipulator~\citep{sk2024boltsorting}, a direction the fifth study's joint bolt/nut detection and
    the eighth study's peg-in-hole pre-assembly alignment later advance, though neither closes it as
    the original study specifically framed it.
\end{description}

%% ------------------------------------------------------------
\subsection{G2: Runtime Arm Selection from Onboard Perception}\label{sec:concl_g2}
%% ------------------------------------------------------------

\begin{description}
    \item[Contribution.] The third study lets two Kinova Gen3 Lite manipulators, each with its own
    wrist-mounted eye-in-hand camera, independently search a shared workspace and relay a detected
    target's pose across a calibrated inter-arm transform, resolving which arm should act through a
    runtime, reachability-based role-assignment rule, validated end-to-end in Gazebo/ROS2 and on
    hardware with sub-5\,mm, sub-1.3$^\circ$ manipulation accuracy regardless of which arm is
    assigned~\citep{sk2025sharedvision}.

    \item[Limitations.] The study's own scope statement notes that its five trials per arm-selection
    configuration establish repeatability at a proof-of-concept level rather than a fully powered
    statistical characterization; that the reachability test is a simple axis-aligned workspace box
    rather than a full manipulability or collision-aware reachability model; that the tie-break used
    when a target lies in both arms' reachable workspaces is plain Euclidean distance rather than any
    measure of manipulability or motion cost; and that role assignment is decided once, at the moment
    of detection, the system does not monitor whether the acting arm's own view of the target remains
    usable once manipulation begins~\citep{sk2025sharedvision}.

    \item[Future directions.] The study identifies extending the protocol so that the idle arm
    continues observing the target during the acting arm's approach, relaying a corrected pose if the
    acting arm's own eye-in-hand view degrades mid-task, as its most direct avenue for future work,
    moving the coordination decision from a one-time handoff to a continuously monitored one. Beyond
    this, it identifies testing the protocol against a larger and more varied set of trials to
    characterize failure modes more rigorously, replacing the fixed workspace-box reachability test
    with a learned or manipulability-weighted reachability measure, and extending the single-object
    handling demonstrated here to multi-object scheduling across the two
    arms~\citep{sk2025sharedvision}.
\end{description}

%% ------------------------------------------------------------
\subsection{G3: Sensor-Free, Real-Time Inter-Arm Distance}\label{sec:concl_g3}
%% ------------------------------------------------------------

\begin{description}
    \item[Contribution.] The fourth study models manipulator links as line segments in 3-D space,
    classifies their configuration as parallel, intersecting, or skew, and computes a continuous,
    sensor-free inter-arm distance signal across sixteen collision-critical link pairs in real time
    from forward kinematics and joint feedback alone, validated on dual Kinova Gen3 Lite manipulators
    during twist control and integrated with a traditional artificial potential field controller to
    demonstrate the signal's direct usability for collision
    avoidance~\citep{sk2025distance}.

    \item[Limitations.] The study states its framework assumes the shared workspace is free of
    external obstacles, and that its practical validation is limited to dual-arm
    systems~\citep{sk2025distance}.

    \item[Future directions.] The study identifies extending the framework to multi-arm systems with
    more than two arms and with variable degrees of freedom, and advancing collision-avoidance
    techniques that consume the real-time distance signal for complex coordinated motions in
    multi-arm systems, as its stated directions for future
    work~\citep{sk2025distance}.
\end{description}

%% ------------------------------------------------------------
\subsection{G4: Stable Asymmetric Dual-Arm Motion Planning}\label{sec:concl_g4}
%% ------------------------------------------------------------

\begin{description}
    \item[Contribution.] The fifth study proposes iAPF, an improved artificial potential field
    framework coupling exponential goal attraction, continuous inter-arm repulsion, and home-seeking
    forces in a three-way equilibrium stabilized by velocity-dependent damping, combined with an
    empirically validated, stable hybrid twist controller and a priority-lock coordination
    mechanism, driven by YOLOv8 OBB detection at 20 frames per second with 0.99 and 0.97 accuracy
    for bolts and nuts respectively~\citep{surya2025iapf}.

    \item[Limitations.] The study states that its framework assumes a controlled workspace free from
    external dynamic obstacles, focusing primarily on inter-arm collision avoidance rather than
    general obstacle avoidance, and that its vision system is currently limited to the pre-trained
    component classes it was trained on, nuts and bolts, though it remains adaptable to new object
    types through retraining~\citep{surya2025iapf}.

    \item[Future directions.] The study identifies extending the present framework to handle external
    dynamic environments through real-time external obstacle detection using RGB-D depth sensing to
    detect unknown objects in the workspace, and adaptive repulsive force field generation that
    incorporates such external obstacles into the existing iAPF framework, enabling deployment in
    unstructured manufacturing environments while retaining the stability and collision-avoidance
    properties demonstrated in the study. It further identifies automated assembly of industrial
    components as a potential application, transitioning from the absolute motion control
    demonstrated here to relative motion control during an assembly phase, with integrated size
    matching and precise alignment capabilities~\citep{surya2025iapf}, a direction the eighth study's
    peg-in-hole pre-assembly alignment work subsequently advances, though by a learned rather than a
    potential-field controller.
\end{description}

%% ------------------------------------------------------------
\subsection{G5: Adaptive Task Allocation}\label{sec:concl_g5}
%% ------------------------------------------------------------

\begin{description}
    \item[Contribution.] The sixth study replaces the fixed nut-left/bolt-right object assignment of
    the fifth study with a distance-based competitive task allocation mechanism, in which each arm
    evaluates object detections over a five-frame temporal proposal window before committing to a
    runtime, distance-based selection, giving an approximately 47\% velocity improvement over fixed
    allocation in homogeneous scenarios with zero collisions across all trials, and universal
    object-handling capability in mixed scenarios at an approximately 11\% time
    overhead~\citep{sk2025competitive}.

    \item[Limitations.] The study's own future-work framing indicates that the demonstrated task is
    absolute-position pick-and-sort placement rather than assembly; that the temporal window length
    and the improved artificial potential field's force coefficients are fixed, hand-set parameters
    rather than ones adapted to real-time workspace density or task complexity; and that validation
    covers single-object pick assignment rather than multi-object assembly sequences of varying
    geometric complexity~\citep{sk2025competitive}.

    \item[Future directions.] The study identifies three directions for future work: extending the
    framework to size-based assembly operations, specifically nut-bolt pairings, which necessitates a
    transition from the absolute positioning movements demonstrated here to coordinated, relative
    end-effector movements for precise threaded fastener insertion and bilateral manipulation;
    implementing adaptive learning mechanisms for the dynamic adjustment of temporal window
    parameters and artificial-potential-field force coefficients based on real-time workspace density
    and task complexity; and expanding validation to multi-object assembly sequences of varying
    geometric complexity to further demonstrate industrial scalability and generalization
    capability~\citep{sk2025competitive}.
\end{description}

%% ------------------------------------------------------------
\subsection{G6: Trajectory Tracking on Non-Spherical Wrists}\label{sec:concl_g6}
%% ------------------------------------------------------------

\begin{description}
    \item[Contribution.] The seventh study presents a coupled-reward Proximal Policy Optimization
    controller that maps instantaneous Cartesian pose error directly to joint velocity commands for
    the Kinova Gen3 Lite's non-spherical wrist, trained purely on 1.5 million static reachable poses
    with no trajectory supervision, generalizing zero-shot to circular, square, lemniscate, and 3-D
    Lissajous trajectories on physical hardware, and, on the circle, the only method among the four
    deployed whose position error did not increase from simulation to hardware~\citep{sk2025diktracking}.

    \item[Limitations.] The study states that hardware validation is restricted to a single
    manipulator, the Kinova Gen3 Lite, so generalization to other non-spherical-wrist platforms such
    as KUKA, ABB, or UR variants with offset wrist axes is untested and the cost of porting the
    policy to a different kinematic structure is unknown; that the TRPO-D and SAC-S baselines and the
    reward-ablation variants were trained with a single seed, and all physical experiments are
    single-trial executions, so hardware run-to-run variance is uncharacterized; that the speed
    evaluation is restricted to simulation and the empirical evaluation to four trajectory
    geometries at a single spatial scale; that hardware baseline comparison against TRAC-IK, DLS, and SAC-S was restricted to
    the circular trajectory, while the proposed method was further validated on the lemniscate,
    3-D Lissajous, and square trajectories without retraining, leaving a head-to-head hardware
    comparison on the more complex trajectories open; that controller behaviour near kinematic
    singularities remains uncharacterized, since the policy was not explicitly trained on
    singularity-proximal configurations; and that, because the policy has no explicit dynamic model
    or closed-form stability guarantee, it does not offer the formal safety guarantees available to
    model-based controllers~\citep{sk2025diktracking}.

    \item[Future directions.] The study identifies incorporating singularity-proximal configurations
    during rollout collection to enable verified recovery behaviour at kinematic boundaries;
    multi-seed training of the TRPO-D and SAC-S baselines and the ablation variants; adding a reference
    velocity or waypoint preview to the observation; domain randomization and an external safety
    layer; extending the hardware baseline comparison to the non-circular
    trajectories; and validating on additional non-spherical-wrist manipulators to test whether the
    consistency observed here is specific to the Kinova Gen3 Lite or a more general property of the
    reward design. It further identifies extending the framework to dual-arm manipulation, where the
    learned differential inverse-kinematics controller would serve as a per-arm low-level policy for
    coordinated bi-manual control, as a natural extension toward collaborative assembly operations
    requiring simultaneous precise trajectory tracking under contact
    constraints~\citep{sk2025diktracking}. The eighth study of this thesis takes up exactly this
    direction, reusing the trained policy, frozen, as the single-arm prior inside GPC.
\end{description}

%% ------------------------------------------------------------
\subsection{G7: Sample-Efficient Dual-Arm Coordination Learning}\label{sec:concl_g7}
%% ------------------------------------------------------------

\begin{description}
    \item[Contribution.] The eighth study reuses the seventh study's trained policy, frozen, as the
    single-arm prior inside GPC (Gated Prior Composition), a hierarchical architecture in which a
    Multi-Agent PPO coordination policy, trained under Centralized Training with Decentralized
    Execution, is combined with the frozen prior through a learned, incrementally updated gating
    scalar co-trained with it from timestep zero, achieving 95.9\% reach and 93.9\% alignment success
    (mean over three training seeds) across 4096 parallel simulation environments and 93.3\% reach with 90.0\% alignment success over
    30 physical hardware trials of cooperative peg-in-hole pre-assembly alignment, with policy
    inference at 200\,Hz on a CPU-only platform against the hardware's 20\,Hz control
    loop~\citep{sk2025alpha}.

    \item[Limitations.] The study's discussion states that the Kinova Gen3 Lite's non-spherical,
    force-sensor-less gripper limits grasp verification to proximity thresholds, so any in-hand slip
    during insertion propagates as an unobservable error into alignment; that object pose is assumed
    known a priori, an intentional scope boundary the study describes as leaving the system not yet
    perception-complete rather than an architectural limitation of the policy itself; that the
    simulation's proximity-triggered snap mechanism was excluded from hardware deployment and
    replaced with staged spatial waypoints, a deployment modification the study attributes to the
    mismatch between simulated contact observations and hardware; and that hardware evaluation over
    30 trials yields correspondingly wide 95\% Wilson confidence intervals, 78.7\% to 98.2\% for
    reach and 74.4\% to 96.5\% for alignment~\citep{sk2025alpha}.

    \item[Future directions.] The study identifies integrating force sensing for compliant contact
    handling, which it notes would directly address the grasp-verification limitation without
    requiring structural changes to the blending or reward design, since force observations could be
    appended directly to the local and global critic states already used; integrating visual
    perception for object pose estimation, which would make the system perception-complete without
    modifying the learned policy or reward structure; and extending the framework to $N>2$ arm
    configurations, which the study notes is architecturally supported by the Dec-POMDP formulation
    already used, extension requiring only an expansion of the global state dimensionality and the
    addition of agents to the MAPPO update, with no structural change to the blending mechanism or
    reward design~\citep{sk2025alpha}.
\end{description}

%% ============================================================
\section{Concluding Remarks}\label{sec:concl_remarks}
%% ============================================================

Read together, the future directions identified independently across these eight studies converge
on a small number of recurring themes rather than eight unrelated wish-lists. Perception
completeness recurs three times: the first study's own call to integrate with a dual-arm system, the
second study's call to handle occluded and moving objects, and the eighth study's call to give GPC
a vision front-end so that object pose is sensed rather than assumed, are each, in the terms their
own authors used, instances of the same unfinished step. External-obstacle awareness recurs twice,
in the fourth study's assumption of an obstacle-free shared workspace and the fifth study's identical
assumption for the iAPF controller built on top of it, both flagged as the immediate next extension
by their own authors. Assembly, rather than sorting, recurs three times, in the second study's wish
to extend to nut-bolt pairs, the fifth study's proposed transition from absolute to relative motion
control during an assembly phase, and the sixth study's identically framed call for size-based
nut-bolt assembly, of which the eighth study's peg-in-hole pre-assembly alignment is a first,
partial, and independently arrived-at instance. And the seventh study's own proposed extension to
dual-arm manipulation, with the learned controller reused as a per-arm low-level policy, is precisely
what the eighth study carries out. None of this is a plan this thesis sets out in advance; it is a
convergence visible only in hindsight, across eight studies conducted as a single, connected
progression from single-arm classical vision to dual-arm reinforcement learning, each addressing a
research gap the previous study's classical or learned components left open, and each still leaving,
in its own authors' words, a next gap of exactly the kind the following study goes on to close.
